% This is part of Mes notes de mathématique
% Copyright (c) 2011-2020
%   Laurent Claessens
% See the file fdl-1.3.txt for copying conditions.

%+++++++++++++++++++++++++++++++++++++++++++++++++++++++++++++++++++++++++++++++++++++++++++++++++++++++++++++++++++++++++++
\section{Groupes abéliens finis}
%+++++++++++++++++++++++++++++++++++++++++++++++++++++++++++++++++++++++++++++++++++++++++++++++++++++++++++++++++++++++++++

Source : \cite{FabricegPSFinis}.

\begin{definition}  \label{DefvtSAyb}
    L'\defe{exposant}{exposant!d'un groupe} du groupe \( G\) est le plus petit entier non nul \( n\) tel que \( g^n=e\) pour tout \( g\in G\). S'il n'existe pas, nous disons que l'exposant du groupe est infini.
\end{definition}

\begin{proposition}\label{PROPooSWHHooOzqWkw}
    Si l'ensemble des ordres de tous les éléments d'un groupe est majoré, alors l'exposant du groupe est le plus petit commun multiple des ordres des éléments du groupe

    Pour un groupe fini, l'exposant est le $\ppcm$ des ordres des éléments du groupe.
\end{proposition}

Le théorème de Burnside~\ref{ThooJLTit} nous donnera un bon paquet d'exemples de groupes d'exposant fini dans \( \GL(n,\eC)\).


Dans le cas des groupes abéliens finis, l'exposant joue un rôle important du fait qu'il existe un élément dont l'ordre est l'exposant. Cela est le théorème suivant.

\begin{theorem}[Exposant dans un groupe abélien fini]
    Un groupe abélien fini contient un élément dont l'ordre est l'exposant du groupe.
\end{theorem}

\begin{proof}
    Soit \( G\) un groupe abélien fini et \( x\in G\), un élément d'ordre maximum \( m\). Nous montrons par l'absurde que l'ordre de tous les éléments de \( G\) divise \( m\). Soit donc \( y\in G\), un élément dont l'ordre ne divise pas \( m\); nous notons $q$ son ordre. Vu que \( q\) ne divise pas \( m\), le nombre \( q\) possède au moins un facteur premier plus de fois que \( m\) : soit \( p\) premier tel que la décomposition de \( q\) contienne \( p^{\beta}\) et celle de \( m\) contienne \( p^{\alpha}\) avec \( \beta>\alpha\). Autrement dit,
    \begin{subequations}
        \begin{align}
            m=p^{\alpha}m'\\
            q=p^{\beta}q'
        \end{align}
    \end{subequations}
    où \( m'\) et \( q'\) ne contiennent plus le facteur \( p\). L'élément \( x\) étant d'ordre \( m\), l'élément \( x^{p^{\alpha}}\) est d'ordre \( m'\). De la même manière, l'élément \( y^{q'}\) est d'ordre \( p^{\beta}\). Étant donné que \( p^{\beta}\) et \( m'\) sont premiers entre eux, l'élément  \( x^{p^{\alpha}}y^{q'}\) est d'ordre \( p^{\alpha}m'>m\). D'où une contradiction avec le fait que \( x\) était d'ordre maximal.

    Par conséquent l'ordre de tous les éléments de $G$ divise celui de \( x\) qui est alors le \( \ppcm\) des ordres de tous les éléments de \( G\), c'est-à-dire l'exposant de \( G\).
\end{proof}

\begin{proposition} \label{PropfPRVxi}
    Soit \( G\) un groupe abélien fini et \( x\in G\), un élément d'ordre maximum. Alors
    \begin{enumerate}
        \item
            Il existe un morphisme \( \varphi\colon G\to \gr(x)\) tel que \( \varphi(x)=x\).
        \item   \label{ItemKRYwjU}
            Il existe un sous-groupe \( K\) de \( G\) tel que \( G=\gr(x)\oplus K\).
    \end{enumerate}
\end{proposition}

\begin{proof}
    Nous notons \( a\) l'ordre de \( x\) qui est également l'exposant du groupe \( G\).

    Nous allons prouver la première partie par récurrence sur l'ordre du groupe. Si \( G=\gr(x)\), alors c'est évident. Soit \( H\) un sous-groupe propre de \( G\) contenant \( x\) et tel que le problème soit déjà résolu pour \( H\) : il existe un morphisme \( \varphi\colon H\to \gr(x)\) tel que \( \varphi(x)=x\). Soit \( y\in G\setminus H\), d'ordre \( b\). Nous allons trouver un morphisme $\hat\varphi\colon \gr(H,y)\to \gr(x) $ telle que \( \hat\varphi(x)=x\).

    Pour cela nous commençons par construire les applications suivantes :
    \begin{equation}
        \begin{aligned}
            \tilde \varphi\colon \eZ/b\eZ\times H&\to \gr(x) \\
            (\bar k,h)&\mapsto x^{kl}\varphi(h)
        \end{aligned}
    \end{equation}
    où \( l\) est encore à déterminer, et
    \begin{equation}
        \begin{aligned}
            p\colon \eZ/b\eZ\times H&\to \gr(y,H) \\
            (\bar k,h)&\mapsto y^kh.
        \end{aligned}
    \end{equation}
    Pour que \( \tilde \varphi\) soit bien définie, il faut que \( a\) divise \( bl\). L'application \( p\) est bien définie parce que \( \bar k\) est pris dans \( \eZ/b\eZ\) et que \( b\) est l'ordre de \( y\).

    Nous allons construire le morphisme \( \hat \varphi\) en considérant le diagramme
    \begin{equation}
    \xymatrix{%
    \ker(p) \ar@{^{(}->}[r]        &   \eZ/b\eZ\times H\ar[d]_{\tilde \varphi}\ar[r]^p&\gr(y,H)\ar[ld]^{\hat \varphi}\\
          &   \gr(x)
       }
    \end{equation}
    que l'on voudra être commutatif. Vu que \( p\) est surjective, les théorèmes d'isomorphismes nous disent que
    \begin{equation}
        \gr(y,H)\simeq\frac{ \eZ/b\eZ\times H }{ \ker p }.
    \end{equation}
    Si \( [\bar k,h]\) est la classe de \( (\bar k,h)\) modulo \( \ker(p)\) alors nous voudrions définir \( \hat \varphi\) par
    \begin{equation}        \label{EqeesVxc}
        \hat\varphi\big( [\bar k,h] \big)=\tilde \varphi(\bar k,h).
    \end{equation}
    Pour que cela soit bien définit, il faut que si \( (\bar r,z)\in \ker p\),
    \begin{equation}
        \hat\varphi\big( [\bar k\bar r,hz] \big)=\hat\varphi\big( [\bar k,h] \big),
    \end{equation}
    c'est-à-dire que \( \tilde \varphi(\bar r,z)=e\). Du coup la définition \eqref{EqeesVxc} n'est bonne que si et seulement si
    \begin{equation}
        \ker(p)\subset\ker(\tilde\varphi ).
    \end{equation}
    Nous pouvons obtenir cela en choisissant bien \( l\).

    Déterminons d'abord le noyau de \( p\). Pour cela nous considérons un nombre \( \beta\) divisant \( b\) tel que \( \gr(y)\cap H=\gr(y^{\beta})\). Nous aurons \( p(\bar k,h)=e\) si et seulement si \( y^h=e\). En particulier \( h=y^{-k}\in\gr(y)\cap H=\gr(y^{\beta})\). Si \( h=(y^{\beta})^m=y^{m\beta}\), alors \( k=-m\beta\) et nous avons
    \begin{equation}
        \ker(p)=\{ (-m\beta,y^{m\beta})\tq m\in \eZ \}.
    \end{equation}
    En plus court : \( \ker(p)=\gr(\beta,y^{-\beta})\). Nous devons donc fixer \( l\) de telle sorte que \( \tilde \varphi(\beta,y^{-\beta})=e\). Étant donné que \( \varphi\) prend ses valeurs dans \( \gr(x)\), il existe un entier \( \alpha\) tel que \( \varphi(y^{-\beta})=x^{\alpha}\); en utilisant cet \( \alpha\), nous écrivons
    \begin{equation}
        \tilde \varphi(\beta,y^{-\beta})=x^{\beta l}\varphi(y^{-\beta})=x^{\beta l+\alpha}.
    \end{equation}
    Par conséquent nous choisissons \( l=-\alpha/\beta\). Nous devons maintenant vérifier que ce choix est légitime, c'est-à-dire que \( a\) divise \( bl\) et que \( \alpha/\beta\) est un entier.

    Étant donné que \( y\) est d'ordre \( b\),
    \begin{equation}
        e=\varphi(y^b)=\varphi(y^{-\beta b/\beta})=\varphi(y^{-\beta})^{b/\beta}=x^{b\beta/\alpha}.
    \end{equation}
    Par conséquent \( a\) divise \( \frac{ b\alpha }{ \beta }=-bl\).

    Pour voir que \( l\) est entier, nous nous rappelons que \( a\) est l'exposant de \( G\) (parce que \( x\) est d'ordre maximum) et que par conséquent \( b\) divise \( a\). Mais \( a\) divise \( \alpha\frac{ b }{ \beta }\). Donc \( \alpha/\beta\) est entier.

    Nous passons maintenant à la seconde partie de la preuve. Nous considérons un morphisme \( \varphi\colon G\to \gr(x)\) tel que \( \varphi(x)=x\). La première partie nous en assure l'existence. Nous montrons que
    \begin{equation}
        \begin{aligned}
            \psi\colon G&\to \gr(x)\oplus \ker(\varphi) \\
            g&\mapsto \big( \varphi(g),g\varphi(g)^{-1} \big)
        \end{aligned}
    \end{equation}
    est un isomorphisme. D'abord \( g\varphi(g)^{-1}\) est dans le noyau de \( \varphi\) parce que \( \varphi(g)^{-1}\) étant dans \( \gr(x)\), et \( \varphi\) étant un morphisme,
    \begin{equation}
        \varphi\big( g\varphi(g)^{-1} \big)=\varphi(g)\varphi(g)^{-1}=e.
    \end{equation}
    L'application \( \psi\) est un morphisme parce que, en utilisant le fait que \( G\) est abélien,
    \begin{subequations}
        \begin{align}
            \psi(g_1g_2)&=\big( \varphi(g_1g_2),g_1g_2\varphi(g_1g_2)^{-1} \big)\\
            &=\big( \varphi(g_1)\varphi(g_2),g_1\varphi(g_1)^{-1}g_2\varphi(g_2)^{-1} \big)\\
            &=\psi(g_1)\psi(g_2).
        \end{align}
    \end{subequations}
    L'application \( \psi\) est injective parce que si \( \psi(g)=(e,e)\) alors \( \varphi(g)=e\) et \( g\varphi(g)^{-1}=e\), ce qui implique \( g=e\).

    Enfin \( \psi\) est surjective parce qu'elle est injective et que les ensembles de départ et d'arrivée ont même cardinal. En effet par le premier théorème d'isomorphisme (théorème~\ref{ThoPremierthoisomo}) appliqué à \( \varphi\) nous avons
    \begin{equation}
        | G |=| \gr(x) |\cdot | \ker(\varphi) |.
    \end{equation}
\end{proof}

\begin{theorem} \label{ThoRJWVJd}
    Tout groupe abélien fini (non trivial) se décompose en
    \begin{equation}
        G\simeq \eZ/d_1\eZ\oplus\ldots\oplus \eZ/d_r\eZ
    \end{equation}
    avec \( d_1\geq 1\) et \( d_i\) divise \( d_{i+1}\) pour tout \( i=1,\ldots, r-1\).

    De plus la liste \( (d_1,\ldots, d_r)\) vérifiant ces propriétés est unique.
\end{theorem}

\begin{proof}
    Soit \( x_1\) un élément d'ordre maximal dans \( G\). Soit \( n_1\) son ordre et
    \begin{equation}
        H_1=\gr(x_1)=\eF_{n_1}.
    \end{equation}
    D'après la proposition~\ref{PropfPRVxi}\ref{ItemKRYwjU}, il existe un supplémentaire \( K_1\) tel que \( G=\eF_{n_1}\oplus K_1\). Si \( K_1=\{ 1 \}\) on s'arrête et on garde \( G=\eF_{n_1}\). Sinon on continue de la sorte en prenant \( x_2\) d'ordre maximal dans \( K_1\) etc.

    Nous devons maintenant prouver l'unicité de cette décomposition. Soit
    \begin{equation}
        G=\eF_{d_1}\oplus\ldots\oplus \eF_{d_r}=\eF_{s_1}\oplus\ldots\oplus \eF_{s_q}.
    \end{equation}
    L'exposant de \( G\) est \( d_r\) et \( s_q\). Donc \( d_r=s_q\). Les complémentaires étant égaux nous avons
    \begin{equation}
        \eF_{d_1}\oplus\ldots\oplus \eF_{d_{r-1}}=\eF_{s_1}\oplus\ldots\oplus \eF_{s_{q-1}}.
    \end{equation}
    En continuant nous trouvons \( r=q\) et \( d_i=s_i\).
\end{proof}

%+++++++++++++++++++++++++++++++++++++++++++++++++++++++++++++++++++++++++++++++++++++++++++++++++++++++++++++++++++++++++++
\section{Groupe monogène}
%+++++++++++++++++++++++++++++++++++++++++++++++++++++++++++++++++++++++++++++++++++++++++++++++++++++++++++++++++++++++++++
\label{SECooXIHPooWVSjhT}

Le théorème suivant donne quelques informations à propos de groupes monogènes. Il impliquera dans le corollaire~\ref{CORooMBLSooMHKmAq} qu'un groupe monogène d'ordre \( n\) possède \( \varphi(n)\) générateur où \( \varphi\) est la fonction indicatrice d'Euler définie en~\ref{DEFooWYIGooRVBTil}.

\begin{theorem}     \label{THOooDOMZooOEYHAe}
    Un groupe monogène est abélien. Plus précisément,
    \begin{enumerate}
        \item
            un groupe monogène infini est isomorphe à \( \eZ\),
        \item
            un groupe monogène fini est isomorphe à \( \eZ/n\eZ\) pour un certain \( n\).
    \end{enumerate}
\end{theorem}

\begin{proof}
    Le groupe est abélien parce que $g=a^n$, \( g'=a^{n'}\) implique \( gg'=q^{n+n'}=g'g\). Nous considérons un générateur \( a\) de \( G\) (qui existe parce que $G$ est monogène) et le morphisme surjectif
    \begin{equation}
        \begin{aligned}
            f\colon \eZ&\to G \\
            p&\mapsto a^p.
        \end{aligned}
    \end{equation}
    Si \( G\) est infini, alors \( f\) est injective parce que si \( a^n=a^{n'}\), alors \( a^{n-n'}=e\), ce qui rendrait \( G\) cyclique et par conséquent non infini. Nous concluons que si \( G\) est infini, alors \( f\) est une bijection et donc un isomorphisme \( \eZ\simeq G\).

    Si \( G\) est fini, alors \( f\) n'est pas injective et a un noyau \( \ker f\). Étant donné que \( \ker f\) est un sous-groupe de \( G\), il existe un (unique) \( n\) tel que \( \ker f=n\eZ\) et le premier théorème d'isomorphisme (théorème~\ref{ThoPremierthoisomo}) nous indique que
    \begin{equation}
        \eZ/\ker f=\eZ/n\eZ=\Image f=G.
    \end{equation}

\end{proof}

Le lemme suivant donne une démonstration alternative, avec une construction plus explicite de l'isomorphisme.

\begin{lemma}[\cite{MonCerveau}]   \label{LemZhxMit}

    À propos de groupes monogènes\footnote{Définition~\ref{DEFooWMFVooLDqVxR}.}

    \begin{enumerate}
        \item


    Soit un groupe monogène \( G\) d'ordre fini \( n\) dont \( g\) est un générateur. Alors il existe un isomorphisme
    \begin{equation}
        \phi\colon G\to (\eZ/n\eZ,+)
    \end{equation}
    tel que \( \phi(g)=1\).

\item

    Si \( G\) est un groupe monogène d'ordre infini et si \( g\) est un générateur, alors il existe un isomorphisme
    \begin{equation}
        \phi\colon G\to (\eZ,+)
    \end{equation}
    tel que \( \phi(g)=1\).

   \item

    Soient \( G\) et \( H\) deux groupes monogènes de même ordre. Soient \( g\) un générateur de \( G\) et \( h\), un générateur de \( H\). Il existe un isomorphisme de \( G\) sur \( H\) qui envoie \( g\) sur \( h\).
    \end{enumerate}
\end{lemma}

\begin{proof}

    Commençons par enfoncer une porte ouverte : vu que le groupe est monogène, l'ordre du groupe est égal à l'ordre de son générateur. Nous séparons les cas selon quel l'ordre soit fini ou non.

    \begin{subproof}
        \item[L'ordre de \( G\) est fini et vaut \( n\)]
            Si \( k\in\eZ\), nous notons \( [k]_n\) la classe de \( k\) modulo \( n\), c'est-à-dire l'ensemble \( \{ k+pn\tq p\in \eZ \}\).

            Nous construisons l'isomorphisme \( \phi\colon G\to \eZ/n\eZ\) de la façon suivante:
            \begin{equation}
                \phi(g^m)=[m]_n.
            \end{equation}
            Cela est une bonne définition parce que une égalité du type \( g^m=g^{m'}\) implique que \( m\) et \( m'\) soient dans la même classe modulo \( n\). Nous vérifions que cela est une isomorphisme entre \( G\) et \( \eZ/n\eZ\).

            \begin{subproof}
            \item[Morphisme]
                Pour l'identité, si \( x=e\) alors \( m=0\) et \( \phi(e)=[0]_n\). Et si \( x=g^k\), \( y=g^l\) alors \( \phi(xy)=\phi(g^{k+l})=[k+l]_n=[k]_n+[l]_n=\phi(x)+\phi(y) \).
            \item[Injectif]
                Supposons \( \phi(g^k)=\phi(g^l)\) avec \( k\geq l\). Nous avons \( h^k=h^l\), dont \( h^{k-l}=e\), ce qui donne \( k-l\in [0]_n\) ou encore \( [k]_n=[l]_n\). En particulier \( g^k=g^l\).
            \item[Surjectif]
                La classe \( [k]_n\) est l'image de \( g^k\).
            \end{subproof}

        \item[L'ordre de \( G\) est infini]

                Si l'ordre de \( G\) est infini alors un élément \( x\in G\) s'écrit de façon unique sous la forme \( x=g^m\) avec \( m\in \eZ\). Dans ce cas nous définissons directement \( \phi(g^m)=m\).

                Le reste de la preuve est alors identique au cas d'ordre fini, mais sans les complications liées au modulo.

    \end{subproof}

    La dernière assertion s'obtient des précédentes par composition d'isomorphismes.

\end{proof}


\begin{corollary}\label{CORooMBLSooMHKmAq}
    Un groupe monogène d'ordre \( n\) possède \( \varphi(n)\) générateurs où \( \varphi\) est la fonction indicatrice d'Euler définie en~\ref{DEFooWYIGooRVBTil}.
\end{corollary}

\begin{proof}
    Le théorème~\ref{THOooDOMZooOEYHAe} nous dit qu'un groupe monogène d'ordre \( n\) est isomorphe à \( \eZ/n\eZ\). La proposition~\ref{PropZnmuphiGensn} nous indique que \( \eZ/n\eZ\) possède \( \varphi(n)\) générateurs.
\end{proof}

%++++++++++++++++++++++++++++++++++++++++++++++++++++++++++++++++++++++++++++++++++++++++++++++ 
\section{Ordre d'un élément dans un groupe fini}
%++++++++++++++++++++++++++++++++++++++++++++++++++++++++++++++++++++++++++++++++++++++++++++++
\begin{theorem}[Théorème de Cauchy\cite{ooTZHGooEPFstf}]    \label{THOooSUWKooICbzqM}
    Soit un groupe fini d'ordre \( n\). Pour tout diviseur premier \( p\) de \( n\), le groupe \( G\) possède au moins un élément d'ordre \( p\).
\end{theorem}
\index{théorème!Cauchy!groupe}

Le lemme suivant indique que sous hypothèse de commutativité, l'ordre d'un élément est une notion multiplicative.
\begin{lemma}[\cite{rqrNyg}]    \label{LemyETtdy}
    Soit \( G\) un groupe et \( a,b\in G\) tels que \( ab=ba\) d'ordres respectivement \( r\) et \( s\), deux nombres premiers entre eux. Alors l'élément \( ab\) est d'ordre \( rs\).
\end{lemma}

\begin{proof}
    Étant donné que \( (ab)^{rs}=a^{rs}b^{rs}=1\), l'ordre de \( ab\) divise \( rs\). Et vu que \( r\) et \( s\) sont premiers entre eux, l'ordre de \( ab\) s'écrit sous la forme \( r_1s_1\) avec \( r_1\divides r\) et \( s_1\divides s\). Nous avons
    \begin{equation}
        a^{r_1s_1}b^{r_1s_1}=(ab)^{r_1s_1}=1,
    \end{equation}
    que nous élevons à la puissance \( r_2\) où \( r_2\) est définit en posant \(r=r_1r_2\) :
    \begin{equation}
        a^{rs_1}b^{rs_1}=1.
    \end{equation}
    Et comme \( a^{rs_1}=1\), nous concluons que \( b^{rs_1}=1\). Donc \( s\divides rs_1\). Par le lemme de Gauss \ref{LemPRuUrsD}, nous avons en fait \( s\divides s_1\). Vu qu'on a aussi \( s_1\divides s\), nous avons \( s=s_1\).

    Le même type d'argument donne \( r=r_1\), et finalement l'ordre de \( ab\) est \( r_1s_1=rs\).
\end{proof}

\begin{lemma}[\cite{Combes}]    \label{LemSkIOOG}
    Un sous-groupe d'indice \( 2\) est un sous-groupe normal.
\end{lemma}

\begin{proof}
    Si $H$ est un tel sous-groupe d'un groupe $G$, alors $G$ possède exactement deux classes à gauche par rapport à \( H\) (théorème de Lagrange~\ref{ThoLagrange}) et se partitionne donc en deux parties : \( G=H\cup xH\) avec \( x \notin H \). De même pour les classes à droite : \( G=H\cup Hx\). Puisque la classe à droite \( Hx \) n'est pas $H$, on a \( xH = Hx \), et $H$ est normal dans $G$ par la proposition~\ref{propGroupeNormal}.
\end{proof}

\begin{lemma}[\cite{NielsBMorph}]\label{PropubeiGX}
    Soit \( H\), un sous-groupe normal d'indice \( m\) de \( G\). Alors pour tout \( a\in G\) nous avons \( a^m\in H\).
\end{lemma}

\begin{proof}
    Par définition de l'indice, le groupe \( G/H\) est d'ordre \( m\). Donc si \( [a]\in G/H\), nous avons \( [a]^m=[e]\), ce qui signifie \( [a^m]=[e]\), ou encore \( a^m\in H\).
\end{proof}


\begin{proposition}[\cite{NielsBMorph}]     \label{PROPooVWVIooQzuAlA}
    Soit un groupe fini \( G\) et \( H\), un sous-groupe normal d'ordre \( n\) et d'indice \( m\) avec \( m\) et \( n\) premiers entre eux. Alors \( H\) est l'unique sous-groupe de \( G\) à être d'ordre \( n\).
\end{proposition}

\begin{proof}
    Soit \( H'\) un sous-groupe d'ordre \( n\). Si \( h\in H'\) alors \( h^n=1\) et \( h^m\in H\) par le lemme \ref{PropubeiGX}. Étant donné que \( m\) et \( n\) sont premiers entre eux, par le théorème de Bézout~\ref{ThoBuNjam}, il existe \( a,b\in \eZ\) tels que
    \begin{equation}
        am+bn=1.
    \end{equation}
    Du coup \( h=h^1=(h^m)^a(h^n)^b\). En tenant compte du fait que \( h^n=1\) et \( h^m\in H\), nous avons \( h\in H\). Ce que nous venons de prouver est que \( H'\subset H\) et donc que \( H=H'\) parce que \( | H' |=| H |=| G |/m\).
\end{proof}

\begin{normaltext}
    Notons que cette proposition ne dit pas qu'il existe un sous-groupe d'ordre \( n\) et d'indice \( m\). Il dit juste que s'il y en a un et s'il est normal, alors il n'y en a pas d'autres.
\end{normaltext}

\begin{lemma}       \label{LemqAUBYn}
    L'ensemble des ordres\footnote{Définition \ref{DEFooKWBCooMlmpCP}.} d'un groupe commutatif est stable par PPCM\footnote{Définition \ref{DefrYwbct}.}.

    Autrement dit, si \( x\in G\) est d'ordre \( r\) et si \( y\in G\) est d'ordre \( s\), alors il existe un élément d'ordre \( \ppcm(r,s)\).
\end{lemma}

\begin{proof}
    Soit \( m=\ppcm(r,s)\). Afin d'écrire \( m\) sous une forme pratique, nous considérons les décompositions en facteurs premiers de \( r\) et \( s\) :
    \begin{subequations}
        \begin{align}
            r&=\prod_{i=1}^kp_i^{\alpha_i}\\
            s&=\prod_{i=1}^kp_i^{\beta_i}
        \end{align}
    \end{subequations}
    où \( \{ p_i \}_{i=1\ldots k}\) est l'ensemble des nombres premiers arrivant dans les décompositions de \( r\) et de \( s\). Si nous posons
    \begin{subequations}
        \begin{align}
            r'&=\prod_{\substack{i=1\\\alpha_1>\beta_i}}^kp_i^{\alpha_i}\\
            s'&=\prod_{\substack{i=1\\\alpha_i\leq \beta_i}}^kp_i^{\beta_i},
        \end{align}
    \end{subequations}
    alors \( \ppcm(r,s)=r's'\) et \( r'\) et \( s'\) sont premiers entre eux. L'élément \( x^{r/r'}\) est d'ordre \( r'\) et l'élément \( y^{s/s'}\) est d'ordre \( s'\). Maintenant nous utilisons le fait que \( G\) soit commutatif et le lemme~\ref{LemyETtdy} pour conclure que l'ordre de \( x^{r/r'}y^{s/s'}\) est \( r's'=m\).
\end{proof}


%+++++++++++++++++++++++++++++++++++++++++++++++++++++++++++++++++++++++++++++++++++++++++++++++++++++++++++++++++++++++++++
\section{Automorphismes du groupe \texorpdfstring{$ \eZ/n\eZ$}{Z/nZ}}
%+++++++++++++++++++++++++++++++++++++++++++++++++++++++++++++++++++++++++++++++++++++++++++++++++++++++++++++++++++++++++++

Notons que \( \eZ/n\eZ=\eF_n\) est un groupe pour l'addition tandis que \( (\eZ/n\eZ)^*\) est un groupe pour la multiplication. Il ne peut donc pas y avoir d'équivoque.

\begin{theorem}[\cite{ooUDKTooTWYpzN}]   \label{ThoozyeSn}
    Pour chaque \( x\in (\eZ/n\eZ)^*\) nous considérons l'application
    \begin{equation}
        \begin{aligned}
            \sigma_x\colon \eZ/n\eZ&\to \eZ/n\eZ \\
            y&\mapsto xy.
        \end{aligned}
    \end{equation}
    L'application
    \begin{equation}
        \sigma\colon \big( (\eZ/n\eZ)^*,\cdot\big)\to \Aut\big( \eZ/n\eZ,+ \big)
    \end{equation}
    ainsi définie est un isomorphisme de groupes.
\end{theorem}
L'énoncé de ce théorème s'écrit souvent rapidement par
\begin{equation}
    \Aut(\eZ/n\eZ)=(\eZ/n\eZ)^*,
\end{equation}
mais il faut bien garder à l'esprit qu'à gauche on considère le groupe additif et à droite celui multiplicatif.

\begin{proof}
    Nous notons \( [x]\) la classe de \( x\) dans \( \eZ/n\eZ\). Nous avons \( \eZ/n\eZ=[1]\). Soit \( f\) un automorphisme de \( (\eZ/n\eZ,+)\); pour tout \( r\in \eZ\) nous avons
    \begin{equation}
        f([r])=f(r[1])=rf([1])=[r]f([1]).
    \end{equation}
En particulier, vu que \( f\) est surjective, il existe un \( r\) tel que \( f([r])=[1]\). Pour un tel \( r\) nous avons \( [1]=[r]f([1])\), c'est-à-dire que nous avons montré que \( f([1])\) est inversible dans \(  \big( (\eZ/n\eZ)^*,\cdot\big)\). Nous montrons à présent que\footnote{Le \( \sigma\) donné ici est l'inverse de celui donné dans l'énoncé. Cela ne change évidemment rien à la validité de l'énoncé et de la preuve.}
    \begin{equation}
        \begin{aligned}
            \sigma\colon \Aut( (\eZ/n\eZ,+))&\to \big( (\eZ/n\eZ)^*,\cdot \big)\\
            f&\mapsto f([1])
        \end{aligned}
    \end{equation}
    est un isomorphisme.

    Nous commençons par la surjectivité. Soit \( [a]\in (\eZ/n\eZ)^*\). Les élément \( [a]\) et \( [1]\) étant tous deux des générateurs de \( (\eZ/n\eZ,+)\), il existe un automorphisme de \( \eZ/n\eZ\) qui envoie \( [1]\) sur \( [a]\) par le lemme~\ref{LemZhxMit}. Cela prouve la surjectivité de \( \sigma\).

    En ce qui concerne l'injectivité, considérons des automorphismes \( f_1\) et \( f_2\) de \( (\eZ/n\eZ,+)\) tels que \( f_1([1])=f_2([1])\). Les automorphismes \( f_1\) et \( f_2\) prennent la même valeur sur un générateur et donc sur tout le groupe. Donc \( f_1=f_2\).

    Enfin nous prouvons que \( \sigma\) est un morphisme, c'est-à-dire que \( \sigma(f\circ g)=\sigma(f)\sigma(g)\). Nous avons
    \begin{subequations}
        \begin{align}
            f\big( g([1]) \big)&=f\big( g([1])[1] \big)=g([1])f([1])=\sigma(f)\sigma(g).
        \end{align}
    \end{subequations}
\end{proof}

Ce dernier résultat s'étend aux groupes cycliques.
\begin{proposition}     \label{PROPooBZOMooVOHoYf}
    Si \( G\) est un groupe cyclique\footnote{Définition \ref{DefHFJWooFxkzCF}.} d'ordre \( n\), alors
    \begin{equation}
        \Aut(G)=(\eZ/n\eZ)^*.
    \end{equation}
\end{proposition}

\begin{corollary}       \label{CorwgmoTK}
    Si \( p\) divise \( q-1\) alors \( \Aut(\eF_q)\) possède un unique sous-groupe d'ordre \( p\).
\end{corollary}

\begin{proof}
    Si \( a\) est un générateur de \( \eF_q^*\) alors le groupe
    \begin{equation}    \label{EqAdGiil}
        \gr\left( a^{\frac{ q-1 }{ p }} \right)
    \end{equation}
    est un sous-groupe d'ordre \( p\). En ce qui concerne l'unicité, soit \( S\) un sous-groupe d'ordre \( p\). Il est donc d'indice \( (q-1)/p\) dans \( \eF_q^*\) et le lemme~\ref{PropubeiGX} nous enseigne que le groupe donné en \eqref{EqAdGiil} est contenu dans \( S\). Il est donc égal à \( S\) parce qu'il a l'ordre de \( S\). Le fait que \( S\) soit normal est dû au fait que \( \eF_q^*\) est abélien.
\end{proof}

%+++++++++++++++++++++++++++++++++++++++++++++++++++++++++++++++++++++++++++++++++++++++++++++++++++++++++++++++++++++++++++
\section{Permutations, groupe symétrique}
%+++++++++++++++++++++++++++++++++++++++++++++++++++++++++++++++++++++++++++++++++++++++++++++++++++++++++++++++++++++++++++

Nous donnons ici quelques éléments à propos du groupe symétrique. Beaucoup de choses supplémentaires sont reportées à la section \ref{SECooZFYQooFfopMa}. Voir aussi le thème \ref{THEMEooQEEWooXDhvhv}.

\begin{definition}      \label{DEFooJNPIooMuzIXd}
    Soit un ensemble \( E\). Une \defe{permutation}{permutation} de l'ensemble \( E\) est une bijection \( E\to E\). Le \defe{groupe symétrique}{groupe!symétrique} de \( E\) le groupe des bijections \( E\mapsto E\); il est noté \( S_E\).

    Le \defe{groupe symétrique}{groupe!symétrique} \( S_n\)\nomenclature[R]{\( S_n\)}{le groupe symétrique} est le groupe des permutations de l'ensemble \( \{ 1,\ldots,n \}\). C'est donc l'ensemble des bijections \( \{ 1,\ldots, n \}\to\{ 1,\ldots, n \}\).
\end{definition}

\begin{definition}      \label{DEFooXNAFooGTbTTJ}
    Une \defe{transposition}{transposition} est une permutation qui inverse deux éléments. Plus précisément, une bijection \( \sigma\colon E\to E\) est une transposition si il existe \( a,b\in E\) tels que
    \begin{equation}
        \sigma(x)=\begin{cases}
              a  &   \text{si } x=b\\
            b    &   \text{si } x=a\\
            x    &    \text{sinon. }
        \end{cases}
    \end{equation}
\end{definition}

\begin{lemma}[\cite{ooJBHPooToyYYI}]        \label{LEMooSGWKooKFIDyT}
    Le groupe symétrique \( S_n\) est un ensemble fini contenant \( n!\) éléments.
\end{lemma}

\begin{lemma}[\cite{ooJFLYooKMbycW}]        \label{LEMooUPBOooWbwMTx}
    Deux résultats.
    \begin{enumerate}
        \item
            Tout groupe est isomorphe à un sous-groupe d'un groupe symétrique.
        \item
            Tout groupe fini d'ordre \( n\) est isomorphe à un sous-groupe de \( S_n\).
    \end{enumerate}
\end{lemma}

\begin{proof}
    Soit, pour \( g\in G\) donné, l'application
    \begin{equation}
        \begin{aligned}
            \tau_g\colon G&\to G \\
            x&\mapsto gx.
        \end{aligned}
    \end{equation}
    Cela est une bijection de \( G\). En effet, d'une part, \( \tau_g(x)=y\) pour \( x=g^{-1} y\) (surjection) et, d'autre part, \( \tau_g(x)=\tau_g(y)\) implique \( gx=gy\) et donc \( x=y\) (injection).

    Nous avons donc \( \tau_g\in S_G\). De plus l'application
    \begin{equation}
        \begin{aligned}
            \varphi\colon G&\to S_G \\
            g&\mapsto \tau_g
        \end{aligned}
    \end{equation}
    est un morphisme de groupe. Il est injectif parce que si \( \tau_g=\tau_h\) alors \( gx=hx\) pour tout \( x\). En particulier \( g=h\).

    Donc \( \varphi\colon G\to \Image(\varphi)\) est un isomorphisme entre \( G\) et un sous-groupe de \( S_G\).

    Un groupe fini de cardinal \( n\) est isomorphe à un sous-groupe de \( S_G\); or \( S_G\) est isomorphe à un des \( S_n\).
\end{proof}
%---------------------------------------------------------------------------------------------------------------------------
\subsection{Décomposition en cycles}
%---------------------------------------------------------------------------------------------------------------------------

\begin{definition}      \label{DEFooSupportPermutation}
    Le \defe{support}{support!d'une permutation} d'une permutation $\sigma$ est l'ensemble constitué des éléments modifiés par $\sigma$:
    \begin{equation*}
        \supp \sigma = \{ i \in \{1,\ldots,n \} \tq \sigma(i) \neq i\}.
    \end{equation*}
\end{definition}

\begin{definition}  \label{DEFooWPYSooPWuwWO}
    Nous disons qu'un élément \( \sigma \in S_n\) \defe{inverse}{inversion!dans le groupe symétrique} les nombres \( i<j\) si \( \sigma(i)>\sigma(j)\). Soit \( N_\sigma\) le nombre d'inversions que \( \sigma\in S_n\) possède (c'est le nombre de couples \( (i,j)\) avec \( i<j\) tels que \( \sigma(i)>\sigma(j)\)). L'entier
    \begin{equation}
        \epsilon(\sigma)=(-1)^{N_\sigma}
    \end{equation}
    est la \defe{signature}{signature!d'une permutation} de \( \sigma\).
\end{definition}

Un \wikipedia{fr}{Permutation}{élément du groupe symétrique} \( S_n\) peut être décomposé en produit de cycles de supports disjoints de la façon suivante. Pour \( \sigma \in S_n \), nous écrivons d'abord le cycle qui correspond à l'orbite de \( 1\). Ce sera le cycle
\begin{equation}
    (1,\sigma(1),\sigma^2(1),\ldots, \sigma^k(1))
\end{equation}
avec \( \sigma^{k+1}(1)=1\). Ensuite nous recommençons avec le plus petit élément de \( \{ 1,\ldots, n \}\) à ne pas être dans ce cycle, et puis le suivant, etc. La \emph{structure} d'une telle décomposition est la donnée des nombres \( k_i\) donnant le nombre de cycles de longueur \( i\).

\begin{lemma}[\cite{Combes}]        \label{LemmvZFWP}
    Soit \( \sigma=(i_1,\ldots, i_k)\in S_n\), un cycle de longueur \( k\) et \( \theta\in S_n\). Alors
    \begin{equation}
        \theta\sigma\theta^{-1}=\big( \theta(i_1),\ldots, \theta(i_k) \big).
    \end{equation}
    Tous les cycles de longueur \( k\) sont conjugués entre eux.
\end{lemma}

\begin{proposition}[Classes de conjugaison et structure en cycles\cite{UXMTXxl}] \label{PropEAHWXwe}
    Une classe de conjugaison dans \( S_n\) est formée des permutations ayant une décomposition en cycles disjoints de même structure. Autrement dit, deux permutations \( \sigma\) et \( \sigma'\) sont conjuguées si et seulement si le nombre \( k_i\) de cycles de longueur \( i\) dans \( \sigma\) est le même que le nombre \( k'_i\) de cycles de longueur \( i\) dans \( \sigma'\).
\end{proposition}

\begin{proof}
    Soit \( \sigma=c_1\ldots c_m\) la décomposition de \( \sigma\) en cycles de supports disjoints. Les \( c_i\) sont des cycles de supports disjoints. Si \( \tau\) est une permutation, alors
    \begin{equation}
        \sigma'=\tau\sigma\tau^{-1}=(\tau c_1\tau^{-1})\ldots (\tau c_m\tau^{-1}),
    \end{equation}
    mais \( \tau c_i\tau^{-1}\) est un cycle de même longueur que \( c\), puisque le lemme~\ref{LemmvZFWP} nous dit que si \( \sigma=(a_1,\ldots, a_k)\), alors \( \tau c\tau^{-1}=\big( \tau(a_1),\ldots, \tau(a_k) \big)\). Notons encore que les cycles \( \tau c_i\tau^{-1}\) restent à support disjoints.

    Donc tous les éléments de la classe de conjugaison de \( \sigma\) sont des permutations de même structure de \( \sigma\).

    Réciproquement, si \( \sigma'=c'_1\ldots c'_m\) est une décomposition de \( \sigma'\) en cycles disjoints tels que la longueur de \( c_i\) est la même que la longueur de \( c'_i\), alors il suffit de construire des permutations \( \tau_i\) telles que \( \tau_i c_i\tau_i^{-1}=c_i'\), à travers le lemme~\ref{LemmvZFWP}. Comme les supports des $c_i$ et des $c'_i$ sont disjoints, la permutation \( \tau_1\ldots \tau_m\) conjugue \( \sigma\) et \( \sigma'\).
\end{proof}

\begin{example}     \label{EXooQAXRooBsPURs}
    Voyons les classes de conjugaison de \( S_3\). Étant donné que ce groupe agit par définition sur un ensemble à \( 3\) éléments, aucun élément de \( S_3\) ne possède un cycle de plus de \( 3\) éléments. Il y a donc seulement des cycles de longueur deux ou trois (à part les triviaux). Aucun élément de \( S_3\) n'a une décomposition en cycles disjoints contenant deux cycles de deux ou un cycle de deux et un de trois.

    En résumé il y a trois classes de conjugaison dans \( S_3\). La première est celle contenant seulement l'identité. La seconde est celle contenant les cycles de longueur deux et la troisième contient les cycles de longueur \( 3\).

    Ce sont donc
    \begin{subequations}
        \begin{align}
            C_1&=\{ \id \}\\
            C_2&=\{ (1,2),(1,3),(2,3) \}\\
            C_3&=\{ (1,2,3),(2,1,3) \}.
        \end{align}
    \end{subequations}
\end{example}

\begin{example} \label{ExVYZPzub}
    Les classes de conjugaison de \( S_4\). Nous savons que les classes de conjugaison dans \( S_4\) sont caractérisées par la structure des décompositions en cycles (proposition~\ref{PropEAHWXwe}). Le groupe symétrique \( S_4\) possède dont les classes de conjugaison suivantes.
\begin{enumerate}
    \item
        Le cycle vide qui représente la classe constituée de l'identité seule.
    \item
        Les transpositions (de type \( (a,b)\)) qui sont au nombre de \( 6\).
    \item
        Les \( 3\)-cycles. Pour savoir \href{http://www.toujourspret.com/techniques/expression/chants/C/cantique_des_etoiles.php}{quel est leur nombre} nous commençons par remarquer qu'il y a \( 4\) façons de prendre \( 3\) nombres parmi \( 4\) et ensuite \( 2\) façons de les arranger. Il y a donc \( 8\) éléments dans cette classe de conjugaison.
    \item
        Les \( 4\)-cycles. Le premier est arbitraire (parce que c'est cyclique). Pour le second il y a \( 3\) possibilités, et deux possibilités pour le troisième; le quatrième est alors automatique. Cette classe de conjugaison contient donc \( 6\) éléments.
    \item       \label{ITEMooGCMYooKZgFHX}
        Les doubles transpositions, du type \( (a,b)(c,d)\). Dans ce cas, tous les nombres sont permutés, et l'image de $1$ détermine la double transposition. Il y a $3$ images possibles, et donc \( 3\) éléments dans cette classe.
\end{enumerate}
\index{classe de conjugaison!dans $S_4$ }
\end{example}

\begin{proposition} \label{PropPWIJbu}
    Tout élément de \( S_n\) peut être écrit sous la forme d'un produit fini de transpositions.
\end{proposition}
Cette décomposition n'est pas à confondre avec celle en cycles de support disjoints. Par exemple \( (1,2,3)=(1,3)(1,2)\).

\begin{propositionDef}\label{PROPooKRHEooAxtmRv}
    Si une permutation peut être écrire sous forme d'un produit d'un nombre pair de permutations, alors toute décomposition en permutations sera en quantité paire.

    Une telle permutation est une \defe{permutation paire}{permutation paire}.
\end{propositionDef}

\begin{lemma}[\cite{PDFpersoWanadoo}]       \label{LemhxnkMf}
    Un \( k\)-cycle est une permutation impaire si \( k\) est pair et paire si \( k\) est impair.
\end{lemma}

\begin{proposition}[\cite{Combes}]  \label{ProphIuJrC}
    Soit \( S_n\) le groupe symétrique.
    \begin{enumerate}
        \item       \label{ITEMooBQKUooFTkvSu}
            L'application \( \epsilon\colon S_n\to \{ 1,-1 \}\) est l'unique homomorphisme surjectif de \( S_n\) sur \( \{ -1,1 \}\).
        \item
            Si \( s=t_1\cdots t_k\) est le produit de \( k\) transpositions, alors \( \epsilon(s)=(-1)^k\).
    \end{enumerate}
\end{proposition}

\begin{proof}
    Soit \( \sigma,\theta \in S_n\). Afin de montrer que \( \epsilon(\sigma\theta )=\epsilon(\sigma)\epsilon(\theta )\), nous divisons les couples \( (i,j)\) tels que \( i\leq j\) en \( 4\) groupes suivant que \( \theta(i)\gtrless \theta(j)\) et \( \sigma\big( \theta(i) \big)\gtrless \sigma\big( \theta(j) \big)\). Nous notons \( N_1\), \( N_2\), \( N_3\) et \( N_4\) le nombre de couples dans chacun des quatre groupes :
    \begin{center}
    \begin{tabular}{c|c|c}
        $ (i,j)$&   \(\sigma\big( \theta(i) \big)<\sigma\big( \theta(j) \big)\)    &   \(\sigma\big( \theta(i) \big)>\sigma\big( \theta(j) \big)\)\\
        \hline
        \( \theta(i)<\theta(j)\)& \( N_1\)&\( N_2\)\\
        \hline
        \( \theta(i)>\theta(j)\)&\( N_3\)&\( N_4\)
    \end{tabular}
    \end{center}
    Nous avons immédiatement \( N_\theta=N_3+N_4\) et \( N_{\sigma\theta}=N_2+N_4\). Les éléments qui participent à \( N_\sigma\) sont ceux où \( \theta(i)\) et \(\theta(j)\) sont dans l'ordre inverse de \( \sigma\big( \theta(i) \big)\) et \( \sigma\big( \theta(j) \big)\) (parce que \(  \theta\) est une bijection). Donc \( N_\sigma=N_2+N_3\). Par conséquent nous avons
    \begin{equation}
        \epsilon(\sigma)\epsilon(\theta)=(-1)^{N_2+N_3}(-1)^{N_3+N_4}=(-1)^{N_2+N_4}=(-1)^{N_{\sigma\theta}}=\epsilon(\sigma\theta).
    \end{equation}
    Nous avons prouvé que \( \epsilon\) est un homomorphisme. Pour montrer que \( \epsilon\) est surjectif sur \( \{ -1,1 \}\) nous devons trouver un élément \( \tau\in S_n\) tel que \( \epsilon(\tau)=-1\). Si \( \tau\) est la transposition \( 1\leftrightarrow 2\) alors le couple \( (1,2)\) est le seul à être inversé par \( \tau\) et nous avons \( \epsilon(\tau)=-1\).

    Avant de montrer l'unicité, nous montrons que si \( \sigma=t_1\ldots t_k\) alors \( \epsilon(\sigma)=(-1)^k\). Pour cela il faut montrer que \( \epsilon(\tau)=-1\) dès que \( \tau\) est une transposition. Soit \( \tau_{ij}\), la transposition \( (i,j)\) et \( \theta=(i,i+1,\ldots, j-1)\) alors le lemme~\ref{LemmvZFWP} dit que
    \begin{equation}
        \tau_{ij}=\theta\tau_{j-1,j}\theta^{-1}.
    \end{equation}
    La signature étant un homomorphisme,
    \begin{equation}
        \epsilon(\tau_{ij})=\epsilon(\theta)\epsilon(\tau_{j-1,j})\epsilon(\theta)^{-1}=\epsilon(\tau_{j-1,j})=-1.
    \end{equation}

    Nous passons maintenant à la partie unicité de la proposition. Soit un homomorphisme surjectif \( \varphi\colon S_n\to \{ -1,1 \}\) et \( \tau\), une transposition telle que \( \varphi(\tau)=-1\) (qui existe parce que sinon \( \varphi\) ne serait pas surjectif\footnote{Nous utilisons ici le fait que tous les éléments de \( S_n\) sont des produits de transpositions, proposition~\ref{PropPWIJbu}.}). Si \( \tau'\) est une autre transposition, il existe \( \sigma\in S_n\) tel que \( \tau'=\sigma\tau\sigma^{-1}\) (lemme~\ref{LemmvZFWP}). Dans ce cas, \( \varphi(\tau')=\varphi(\tau)=-1\), et si \( \sigma=\tau_1\ldots \tau_k) \),
    \begin{equation}
         \varphi(\sigma)=(-1)^k=\epsilon(\sigma).
    \end{equation}
\end{proof}

\begin{corollary}       \label{CORooZLUKooBOhUPG}
    Si \( \sigma\in S_n\), alors
    \begin{equation}
        \epsilon(\sigma)=\epsilon(\sigma^{-1}).
    \end{equation}
\end{corollary}

\begin{proof}
    Comme dit par la proposition \ref{ProphIuJrC}, \( \epsilon\) est un homomorphisme, donc
    \begin{equation}
        \epsilon(\sigma)\epsilon(\sigma^{-1})=\epsilon(\sigma\sigma^{-1})=\epsilon(\id)=1.
    \end{equation}
    Vu que \( \epsilon(\sigma)\) et \( \epsilon(\sigma^{-1})\) ne peuvent être que \( \pm1\), ils doivent être tous les deux \( 1\) ou tous les deux \( -1\) pour que le produit soit \( 1\).
\end{proof}
